%----------------------------------------------------------------------------------------
% MODELLEREN MET DATA. Alle opgaven zijn eigen werk; de gegevens zijn voorbeeldgegevens.
%
% Opbouw: Start, leerstof (theorie, werkwijze of uitleg, voorbeeld, korte oefening), en op een
% nieuwe pagina de Opdrachten (door elkaar, stijgende moeilijkheid), Koppel en Test jezelf.
% Labels van dit hoofdstuk beginnen met h2-.
%----------------------------------------------------------------------------------------

\mychapter{2}{Modelleren met data}

\lkr{Leerplan B+S'': LPD 55 (onderzoekscompetentie), met het modelleren van meetgegevens. De kleinste-kwadratenrechte wordt hier ingevoerd zoals ze in een onderzoek gebruikt wordt; de formules worden in sectie 2 bewezen met de top van een dalparabool; een afleiding met afgeleiden staat als opdracht. Met een rekentoestel of rekenblad thuis gaat het rekenwerk sneller; in de klas volstaan kleine datasets. Alle getallen in dit hoofdstuk zijn voorbeeldgegevens, geen echte metingen.}

Bij veel onderzoeksvragen meet je iets: hoe lang een slinger is en hoe snel hij slingert, hoe snel thee afkoelt. De punten liggen nooit precies op een mooie grafiek: er zijn meetfouten, en het model is een vereenvoudiging. Toch wil je weten welk verband het best past. Dat is modelleren. Dit hoofdstuk volgt de stappen: variabelen kiezen en meten, een model kiezen, de best passende rechte berekenen, een verband rechttrekken, en het model kritisch bekijken.

\begin{oefeningen}[Start]
\oef{h2-start-1}

\oef{h2-verken}
\end{oefeningen}

%========================================================================================
\section{Meten en een model kiezen}

\subsection{Variabelen}\label{sec:h2-variabelen}

Bij een experiment verander je één grootheid bewust, en kijk je hoe een andere grootheid reageert. Alles wat verder invloed kan hebben, houd je zo goed mogelijk gelijk. Verander je twee dingen tegelijk, dan weet je achteraf niet waardoor een verschil komt.

\begin{theorie}
De \textbf{\textcolor{accent}{onafhankelijke variabele}} is de grootheid die je zelf kiest of instelt. Ze komt op de $x$-as.

\vspace{1mm}
De \textbf{\textcolor{accent}{afhankelijke variabele}} is de grootheid die je meet, en waarvan je verwacht dat ze van de onafhankelijke afhangt. Ze komt op de $y$-as.

\vspace{1mm}
De grootheden die je constant houdt, zijn de \textbf{\textcolor{accent}{controlevariabelen}}.
\end{theorie}

\voorbeeld{}{Je onderzoekt hoe de periode van een slinger (de tijd voor één keer heen en terug) afhangt van zijn lengte. Welke variabelen zijn er?\label{vb:h2-variabelen}}{%
\uitleg{De vraag zegt al wat je kiest en wat je meet. Voor de controlevariabelen zoek je wat de periode nog meer zou kunnen beïnvloeden.}
\begin{stappenvrij}{5.4cm}{9.3cm}
\text{onafhankelijk: lengte } L & Die kies je zelf: bijvoorbeeld $0{,}2$ m, $0{,}4$ m, \ldots, $1{,}0$ m. \\
\text{afhankelijk: periode } T & Die meet je met een chronometer. \\
\text{controle: massa, uitwijking} & Hetzelfde gewichtje, en telkens dezelfde kleine beginuitwijking. Dan kan een verschil in $T$ alleen van $L$ komen.
\end{stappenvrij}}

\begin{oefeningen}
\oef{h2-variabelen}
\end{oefeningen}

\subsection{Meten}

Geen enkele meting is exact. Een chronometer start net te laat, een lat wordt scheef gelegd, een weegschaal staat niet op nul. Die fouten zijn van twee soorten, en je pakt ze anders aan.

\begin{theorie}
De \textbf{\textcolor{accent}{meetfout}} is het verschil tussen de gemeten en de werkelijke waarde. Een \textbf{\textcolor{accent}{toevallige fout}} verschilt van meting tot meting: nu eens te groot, dan weer te klein. Een \textbf{\textcolor{accent}{systematische fout}} zit bij elke meting in dezelfde richting.
\end{theorie}

Toevallige fouten maak je kleiner door te herhalen en het gemiddelde te nemen: de fouten heffen elkaar dan voor een deel op. Een systematische fout verdwijnt niet door te herhalen, want ze zit in elke meting. Die moet je opsporen en wegwerken, bijvoorbeeld door de weegschaal op nul te zetten.

\voorbeeld{}{Je meet de periode $T$ van een slinger van $1{,}00$ m, die ongeveer $2$ s duurt. Door je reactietijd zit je met de chronometer tot $0{,}2$ s naast. Hoe meet je $T$ nauwkeurig?\label{vb:h2-meten}}{%
\uitleg{Een fout van $0{,}2$ s op één periode van $2$ s is $10$\,\%. Die fout maak je kleiner door een langere tijd te meten, en door te herhalen.}
\begin{stappenbreed}
\text{tijd voor 10 periodes: } 20{,}0;\ 20{,}2;\ 20{,}1 \text{ s} & Je reactietijd speelt maar één keer mee, bij start en stop, en niet tien keer. \\
\text{gemiddelde: } 20{,}1 \text{ s} & Drie keer herhalen dempt de toevallige fout. \\
T = \frac{20{,}1}{10} = 2{,}01 \text{ s} & De fout van $0{,}2$ s wordt per periode $0{,}02$ s: nog maar $1$\,\%.
\end{stappenbreed}
Tel je telkens één slingering te weinig, dan maak je een systematische fout: elke meting is dan $10$\,\% te klein, en het gemiddelde ook.}

\begin{oefeningen}
\oef{h2-meten}
\end{oefeningen}

\subsection{Een model kiezen}

Na het meten zet je de punten in een figuur, en stel je een hypothese op: welk soort functie beschrijft het verband? Die hypothese komt uit de theorie (de natuurkunde voorspelt hoe een slinger zich gedraagt), of uit de vorm van de punten.

\begin{theorie}
Een \textbf{\textcolor{accent}{spreidingsdiagram}} toont de metingen als punten $(x_i, y_i)$.

\vspace{1mm}
Een \textbf{\textcolor{accent}{model}} is een functie die het verband tussen de variabelen beschrijft, zoals

\vspace{1mm}
\quad lineair \quad $\formule{y = ax + b}$ \qquad exponentieel \quad $\formule{y = c \cdot g^x}$ \qquad macht \quad $\formule{y = c \cdot x^k}$
\end{theorie}

Zijn de $x$-waarden gelijk verdeeld, dan helpt een tabel bij de keuze. Zijn de verschillen van de $y$-waarden constant, dan is het verband lineair (zie de eigenschap \textit{Constante verschillen} in hoofdstuk 1). Zijn de verhoudingen constant, dan is het exponentieel: uit $y_{i+1} = g \cdot y_i$ volgt op dezelfde manier $y_i = y_0 \cdot g^i$.

\voorbeeld{}{Een biologe telt om het uur het aantal gistcellen $N$ (in miljoen) in een proefbuis: $200$; $240$; $288$; $346$; $415$. Welk model past?\label{vb:h2-model}}{%
\uitleg{De metingen liggen telkens een uur uit elkaar. We berekenen eerst de verschillen, en als die niet constant zijn, de verhoudingen.}
\begin{stappenbreed}
\text{verschillen: } 40;\ 48;\ 58;\ 69 & Ze stijgen: het verband is niet lineair. \\
\text{verhoudingen: } \frac{240}{200} = 1{,}20;\ \ 1{,}20;\ \ 1{,}20;\ \ 1{,}20 & Telkens ongeveer $1{,}20$: elk uur komt er $20$\,\% bij. \\
N = 200 \cdot 1{,}20^t & Een exponentieel model, met $t$ in uur.
\end{stappenbreed}
Dat past ook bij de biologie: elke cel deelt zich, dus hoe meer cellen, hoe meer er bijkomen.}

\begin{oefeningen}
\oef{h2-model}
\end{oefeningen}

%========================================================================================
\section{De best passende rechte}

\subsection{Residuen}

\naast{%
Liggen de metingen $(x_i, y_i)$ ongeveer op een rechte $y = ax + b$, dan zoek je de rechte die zo goed mogelijk bij de punten past. In de Verken legde iedereen de lat net anders. Om te kunnen kiezen, heb je een maat nodig voor hoe ver de punten van de rechte liggen. We kijken naar de verticale afstanden: de fout die de rechte maakt bij elke meting.}{%
\begin{tikzpicture}[x=.7cm,y=.5cm,font=\scriptsize]
  \rooster{0}{1}{5}{0}{1}{6}
  \draw[accent,thick] (0,2.2) -- (5.5,5.5);
  \foreach \x/\y in {1/2,2/4,3/5,4/4,5/5}{
    \draw[grijs,dashed] (\x,\y) -- (\x,{0.6*\x+2.2});
    \fill[accent] (\x,\y) circle (2pt);}
  \node[right] at (3,4.55) {$r_3$};
\end{tikzpicture}}

\begin{theorie}
Bij een rechte $y = ax + b$ is het \textbf{\textcolor{accent}{residu}} van het punt $(x_i, y_i)$ \qquad $\formule{r_i = y_i - (ax_i + b)}$

\vspace{1mm}
De \textbf{\textcolor{accent}{kwadratensom}} van de rechte is \qquad $\formule{S(a, b) = \sum r_i^2 = \sum (y_i - ax_i - b)^2}$

\vspace{1mm}
De \textbf{\textcolor{accent}{kleinste-kwadratenrechte}} is de rechte waarvoor $S(a, b)$ zo klein mogelijk is.
\end{theorie}

Een positief residu betekent dat het punt boven de rechte ligt, een negatief dat het eronder ligt. Waarom kwadraten, en niet gewoon de som van de residuen? Omdat positieve en negatieve residuen elkaar dan opheffen. Zo heeft de horizontale rechte $y = 4$ in de figuur residuen $-2, 0, 1, 0, 1$, met som $0$, terwijl ze de stijging van de punten volledig mist. Kwadraten zijn altijd positief, en ze wegen grote afwijkingen zwaarder dan kleine.

\voorbeeld{}{Een leerling legt bij de punten $(1, 2)$, $(2, 4)$, $(3, 5)$, $(4, 4)$, $(5, 5)$ de lat langs $y = 0{,}5x + 2$. Bereken de residuen en de kwadratensom.\label{vb:h2-residu}}{%
\uitleg{Per punt berekenen we wat de rechte voorspelt, trekken dat af van de meting, en kwadrateren.}
\figuren{\renewcommand{\arraystretch}{1.2}%
\begin{tabular}{|c|c|c|c|c|c|c|}
\hline
$x_i$ & $1$ & $2$ & $3$ & $4$ & $5$ & \\ \hline
$y_i$ & $2$ & $4$ & $5$ & $4$ & $5$ & \\ \hline
$0{,}5x_i + 2$ & $2{,}5$ & $3$ & $3{,}5$ & $4$ & $4{,}5$ & \\ \hline
$r_i$ & $-0{,}5$ & $1$ & $1{,}5$ & $0$ & $0{,}5$ & som $2{,}5$ \\ \hline
$r_i^2$ & $0{,}25$ & $1$ & $2{,}25$ & $0$ & $0{,}25$ & $S = 3{,}75$ \\ \hline
\end{tabular}}
Vier van de vijf residuen zijn positief of nul: de meeste punten liggen boven de lat. De lat ligt dus te laag.}

\begin{oefeningen}
\oef{h2-residu}
\end{oefeningen}

\subsection{De beste rechte met een gegeven richting}

In het voorbeeld lag de lat te laag. Schuif je ze evenwijdig omhoog, dan blijft $a$ gelijk en verandert alleen $b$. Bij een vaste $a$ is $S$ een kwadratische functie van $b$, en waar die het kleinst is, weet je al: in de top van een dalparabool. Zo vind je de beste $b$ zonder te proberen.

\lkr{De formules van de kleinste-kwadratenrechte worden hier bewezen met de top van een dalparabool, in twee stappen. Opdracht \ref{oef:h2-o5} en Koppel-oefening \ref{oef:h2-koppel-1} geven dezelfde afleiding met een tip of met afgeleiden.}

\begin{theorie}
\begin{stelling}{Beste $b$ bij een vaste $a$}
Bij een vaste $a$ is $S(a, b)$ het kleinst voor \quad $\formule{b = \ybar - a\,\xbar}$

\vspace{1mm}
De beste rechte met richtingscoëfficiënt $a$ gaat dus door het punt $\formule{(\xbar, \ybar)}$.
\end{stelling}

\begin{bewijs}
Neem $w_i = y_i - ax_i$. Hun gemiddelde is $\overline{w} = \frac1n \sum (y_i - ax_i) = \ybar - a\,\xbar$, en $\sum w_i = n\,\overline{w}$. Dan is
\bewijsregel{S(a, b) = \sum (w_i - b)^2 = \sum w_i^2 - 2b\sum w_i + nb^2 = nb^2 - 2n\overline{w}\,b + \sum w_i^2 .}
Als functie van $b$ is dat een dalparabool, want $n > 0$. Haar top ligt in
\bewijsregel{b = \frac{2n\overline{w}}{2n} = \overline{w} = \ybar - a\,\xbar .}
Voor die $b$ is $a\,\xbar + b = \ybar$: het punt $(\xbar, \ybar)$ ligt op de rechte.
\end{bewijs}
\end{theorie}

Het punt $(\xbar, \ybar)$ is het zwaartepunt van de punten: het gemiddelde punt. De stelling zegt dat de beste rechte, welke richting je ook kiest, altijd door dat punt gaat. Daar draait de lat als het ware rond.

\voorbeeld{}{Welke rechte met richtingscoëfficiënt $0{,}5$ past het best bij de punten uit voorbeeld~\ref{vb:h2-residu}? Hoeveel kleiner wordt de kwadratensom?\label{vb:h2-besteb}}{%
\uitleg{Volgens de stelling moet de rechte door het gemiddelde punt gaan. Daarna rekenen we de kwadratensom opnieuw uit.}
\begin{stappenbreed}
\xbar = \frac{15}{5} = 3,\quad \ybar = \frac{20}{5} = 4 & Het gemiddelde punt is $(3, 4)$. \\
b = 4 - 0{,}5 \cdot 3 = 2{,}5 & De lat schuift $0{,}5$ omhoog: $y = 0{,}5x + 2{,}5$. \\
r_i\colon\ -1;\ 0{,}5;\ 1;\ -0{,}5;\ 0 & Nu liggen er evenveel punten boven als onder de rechte. \\
S = 1 + 0{,}25 + 1 + 0{,}25 + 0 = 2{,}5 & Kleiner dan $3{,}75$.
\end{stappenbreed}}

\begin{oefeningen}
\oef{h2-besteb}
\end{oefeningen}

\subsection{De richtingscoëfficiënt}

Nu de lat door $(\xbar, \ybar)$ gaat, kunnen we ze nog rond dat punt draaien. Welke richting is de beste? Ook dat blijkt de top van een dalparabool te zijn, nu als functie van $a$.

\begin{theorie}
\begin{stelling}{Kleinste-kwadratenrechte}
Zijn niet alle $x_i$ gelijk, dan is de kleinste-kwadratenrechte $y = ax + b$ met

\vspace{1mm}
\qquad $\formule{a = \frac{\sum (x_i - \xbar)(y_i - \ybar)}{\sum (x_i - \xbar)^2}}$ \qquad en \qquad $\formule{b = \ybar - a\,\xbar}$
\end{stelling}

\begin{bewijs}
Volgens de vorige stelling is $S(a, b) \geqslant S(a, \ybar - a\,\xbar)$ voor elke $a$ en $b$. Met $b = \ybar - a\,\xbar$ is
\bewijsregel{y_i - ax_i - b = (y_i - \ybar) - a(x_i - \xbar) .}
Noteer kort $A = \sum (x_i - \xbar)^2$, $B = \sum (x_i - \xbar)(y_i - \ybar)$ en $C = \sum (y_i - \ybar)^2$. Uitwerken van de kwadraten geeft
\bewijsregel{S(a, \ybar - a\,\xbar) = \sum \bigl((y_i - \ybar) - a(x_i - \xbar)\bigr)^2 = Aa^2 - 2Ba + C .}
Omdat niet alle $x_i$ gelijk zijn, is $A > 0$: een dalparabool in $a$, met top in $a = \frac{2B}{2A} = \frac{B}{A}$. Voor elke $a$ en $b$ is dus
\bewijsregel{S(a, b) \geqslant S(a, \ybar - a\,\xbar) \geqslant S\Bigl(\frac{B}{A},\ \ybar - \frac{B}{A}\,\xbar\Bigr) .}
\end{bewijs}
\end{theorie}

De noemer $A$ meet hoe sterk de $x_i$ rond hun gemiddelde spreiden. De teller $B$ meet of $x$ en $y$ samen boven of samen onder hun gemiddelde liggen: dan is het product positief. Een grote positieve teller betekent dus een stijgende rechte.

\voorbeeld{}{Bepaal de kleinste-kwadratenrechte voor de punten $(1, 2)$, $(2, 4)$, $(3, 5)$, $(4, 4)$, $(5, 5)$.\label{vb:h2-rechte}}{%
\uitleg{We zetten de afwijkingen van het gemiddelde in een tabel en tellen op. Het gemiddelde punt is $(3, 4)$, zoals in voorbeeld~\ref{vb:h2-besteb}.}
\figuren{\renewcommand{\arraystretch}{1.2}%
\begin{tabular}{|c|c|c|c|c|c|c|}
\hline
$x_i - \xbar$ & $-2$ & $-1$ & $0$ & $1$ & $2$ & \\ \hline
$y_i - \ybar$ & $-2$ & $0$ & $1$ & $0$ & $1$ & \\ \hline
$(x_i - \xbar)(y_i - \ybar)$ & $4$ & $0$ & $0$ & $0$ & $2$ & som $6$ \\ \hline
$(x_i - \xbar)^2$ & $4$ & $1$ & $0$ & $1$ & $4$ & som $10$ \\ \hline
\end{tabular}}
\begin{stappenbreed}
a = \frac{6}{10} = 0{,}6 & Teller en noemer uit de laatste kolom. \\
b = 4 - 0{,}6 \cdot 3 = 2{,}2 & De rechte gaat door $(3, 4)$.
\end{stappenbreed}
De rechte is $y = 0{,}6x + 2{,}2$. Ze is iets steiler dan de lat uit voorbeeld~\ref{vb:h2-besteb}, en haar kwadratensom is $2{,}4$: nog iets kleiner dan $2{,}5$.}

\begin{werkwijze}
Bereken $\xbar$ en $\ybar$. Zet in een tabel de afwijkingen $x_i - \xbar$ en $y_i - \ybar$, hun producten en de kwadraten $(x_i - \xbar)^2$. Deel de som van de producten door de som van de kwadraten: dat is $a$. Bereken dan $b = \ybar - a\,\xbar$.
\end{werkwijze}

\begin{oefeningen}
\oef{h2-helling}
\end{oefeningen}

\subsection{De som van de residuen}

In voorbeeld~\ref{vb:h2-residu} lag de lat te laag, en waren de residuen samen $2{,}5$. Voor de beste rechte kan dat niet: daar is de som altijd $0$.

\begin{theorie}
\begin{stelling}{Som van de residuen}
Voor de kleinste-kwadratenrechte is \quad $\formule{\sum r_i = 0}$
\end{stelling}

\begin{bewijs}
Met $\sum y_i = n\,\ybar$, $\sum x_i = n\,\xbar$ en $b = \ybar - a\,\xbar$ is
\bewijsregel{\sum r_i = \sum (y_i - ax_i - b) = n\,\ybar - a\,n\,\xbar - n(\ybar - a\,\xbar) = 0 .}
\end{bewijs}
\end{theorie}

De positieve en negatieve residuen houden elkaar dus precies in evenwicht: de rechte loopt midden door de puntenwolk. Het omgekeerde geldt niet. Ook de horizontale rechte $y = \ybar$ heeft residuen met som $0$, en die is meestal niet de beste. De eigenschap kan dus alleen een fout opsporen: is de som niet $0$, dan heb je zeker niet de kleinste-kwadratenrechte.

\voorbeeld{}{Een leerling vindt voor de punten $(1, 2)$, $(2, 4)$, $(3, 5)$, $(4, 4)$, $(5, 5)$ de rechte $y = 0{,}7x + 2$. Kan dat de kleinste-kwadratenrechte zijn?\label{vb:h2-controle}}{%
\uitleg{We hoeven de formules niet opnieuw te gebruiken: de som van de residuen volstaat als controle.}
\begin{stappenbreed}
\sum r_i = \sum y_i - 0{,}7\sum x_i - 5 \cdot 2 & De som van de residuen, zonder ze een voor een te berekenen. \\
= 20 - 0{,}7 \cdot 15 - 10 = -0{,}5 & Niet nul. \\
0{,}7 \cdot 3 + 2 = 4{,}1 \neq 4 & Hetzelfde zie je aan het gemiddelde punt: $(3, 4)$ ligt niet op de rechte.
\end{stappenbreed}
Dit is dus niet de kleinste-kwadratenrechte. Die is $y = 0{,}6x + 2{,}2$, uit voorbeeld~\ref{vb:h2-rechte}.}

\begin{oefeningen}
\oef{h2-controle}
\end{oefeningen}

\subsection{De correlatiecoëfficiënt}

De kleinste-kwadratenrechte bestaat altijd, ook als de punten helemaal niet op een rechte liggen. Hoe goed past ze? Kijk daarvoor naar de teller $\sum (x_i - \xbar)(y_i - \ybar)$.

\naast{%
Verdeel het vlak met de lijnen $x = \xbar$ en $y = \ybar$ in vier stukken. Rechtsboven en linksonder liggen $x_i$ en $y_i$ samen boven of samen onder hun gemiddelde, en is het product positief. In de andere twee stukken is het negatief. Liggen de punten vooral rechtsboven en linksonder, dan is de som groot en positief. Liggen ze overal verspreid, dan heffen de producten elkaar op. Om een getal te krijgen dat niet van de eenheden afhangt, deel je nog door de spreiding van $x$ en die van $y$.}{%
\begin{tikzpicture}[x=.7cm,y=.5cm,font=\scriptsize]
  \rooster{0}{1}{5}{0}{1}{6}
  \draw[grijs,dashed] (3,0) -- (3,6.2) (0,4) -- (5.3,4);
  \node[text=accent] at (4.6,5.6) {$+$}; \node[text=accent] at (.6,1) {$+$};
  \node[text=grijs] at (.6,5.6) {$-$}; \node[text=grijs] at (4.6,1) {$-$};
  \foreach \x/\y in {1/2,2/4,3/5,4/4,5/5}{\fill[accent] (\x,\y) circle (2pt);}
  \node[below] at (3,-.5) {$\xbar$};
\end{tikzpicture}}

\begin{theorie}
De \textbf{\textcolor{accent}{correlatiecoëfficiënt}} van de punten $(x_i, y_i)$ is

\vspace{1mm}
\qquad $\formule{r = \frac{\sum (x_i - \xbar)(y_i - \ybar)}{\sqrt{\sum (x_i - \xbar)^2 \cdot \sum (y_i - \ybar)^2}}}$
\end{theorie}

De teller van $r$ is dezelfde als die van $a$, en de noemer is positief. Dus $r$ heeft altijd hetzelfde teken als $a$: positief bij een stijgende rechte, negatief bij een dalende.

\voorbeeld{}{Bereken de correlatiecoëfficiënt van de punten $(1, 2)$, $(2, 4)$, $(3, 5)$, $(4, 4)$, $(5, 5)$.\label{vb:h2-correlatie}}{%
\uitleg{De teller en $\sum (x_i - \xbar)^2$ staan al in de tabel van voorbeeld~\ref{vb:h2-rechte}. Alleen $\sum (y_i - \ybar)^2$ is nieuw.}
\begin{stappenbreed}
\sum (y_i - \ybar)^2 = 4 + 0 + 1 + 0 + 1 = 6 & Met de afwijkingen $-2, 0, 1, 0, 1$. \\
r = \frac{6}{\sqrt{10 \cdot 6}} = \frac{6}{\sqrt{60}} \approx 0{,}77 & Positief, zoals $a$.
\end{stappenbreed}
De producten $4, 0, 0, 0, 2$ zijn allemaal positief of nul: in de figuur hierboven ligt geen enkel punt linksboven of rechtsonder. Er is een duidelijk stijgend verband, maar de punten liggen niet mooi op een rechte.}

\begin{oefeningen}
\oef{h2-rechte}

\oef{h2-negatief}
\end{oefeningen}

\subsection{Hoe goed past de rechte?}

Waarom ligt $r$ altijd tussen $-1$ en $1$, en wat betekent een $r$ dicht bij $1$? Het antwoord zit in de kwadratensom van de kleinste-kwadratenrechte. We gebruiken de korte notatie $A$, $B$ en $C$ uit het bewijs van de kleinste-kwadratenrechte, zodat $a = \frac{B}{A}$ en $r = \frac{B}{\sqrt{AC}}$.

\begin{theorie}
\begin{stelling}{Kwadratensom en correlatie}
Zijn niet alle $x_i$ en niet alle $y_i$ gelijk, dan heeft de kleinste-kwadratenrechte als kwadratensom
\bewijsregel{\formule{S_{\min} = (1 - r^2) \cdot \sum (y_i - \ybar)^2}}
Daarom is $\formule{-1 \leqslant r \leqslant 1}$, en is $r = \pm 1$ precies als alle punten op één rechte liggen.
\end{stelling}

\begin{bewijs}
In het bewijs van de kleinste-kwadratenrechte was $S = Aa^2 - 2Ba + C$, met het minimum in $a = \frac{B}{A}$. Invullen geeft
\bewijsregel{S_{\min} = \frac{B^2}{A} - \frac{2B^2}{A} + C = C - \frac{B^2}{A} = C\Bigl(1 - \frac{B^2}{AC}\Bigr) = C(1 - r^2) .}
Een som van kwadraten is nooit negatief, dus $C(1 - r^2) \geqslant 0$. Omdat $C > 0$, is $1 - r^2 \geqslant 0$, of $-1 \leqslant r \leqslant 1$. Verder is $r^2 = 1$ precies als $S_{\min} = 0$, dus als alle residuen $0$ zijn: alle punten liggen op de rechte.
\end{bewijs}
\end{theorie}

De horizontale rechte $y = \ybar$ heeft als kwadratensom $\sum (y_i - \ybar)^2$: dat is de spreiding van de $y_i$ als je $x$ niet gebruikt. De kleinste-kwadratenrechte houdt daarvan een deel $1 - r^2$ over. Bij $r = \pm 1$ blijft er niets over, bij $r = 0$ helpt de rechte niets. Let op: ook $r = 0{,}9$ laat nog $19$\,\% over, want $1 - 0{,}81 = 0{,}19$.

\voorbeeld{}{Bereken de kwadratensom van de kleinste-kwadratenrechte voor de punten van voorbeeld~\ref{vb:h2-correlatie}. Vergelijk met de rechten uit de vorige voorbeelden.\label{vb:h2-rkwadraat}}{%
\uitleg{We hebben $\sum (y_i - \ybar)^2 = 6$ en $r = \frac{6}{\sqrt{60}}$ al berekend. De stelling geeft de kwadratensom zonder één residu uit te rekenen.}
\begin{stappenbreed}
r^2 = \frac{36}{60} = 0{,}6 & Exact, zonder afronden. \\
S_{\min} = (1 - 0{,}6) \cdot 6 = 2{,}4 & De kleinste kwadratensom. \\
6 \ \to\ 3{,}75 \ \to\ 2{,}5 \ \to\ 2{,}4 & Van de rechte $y = 4$ over de lat $y = 0{,}5x + 2$ en $y = 0{,}5x + 2{,}5$ naar $y = 0{,}6x + 2{,}2$.
\end{stappenbreed}
De beste rechte houdt $40$\,\% van de spreiding $6$ over. Met de residuen van $y = 0{,}6x + 2{,}2$ kun je de $2{,}4$ nagaan in opdracht~\ref{oef:h2-o3}.}

\begin{oefeningen}
\oef{h2-rgrenzen}
\end{oefeningen}

%========================================================================================
\section{Niet-lineaire verbanden}

Niet elk verband is lineair. Vaak kun je het wel lineair maken door een van de grootheden, of allebei, te transformeren. Dan bepaal je de kleinste-kwadratenrechte voor de getransformeerde gegevens, en reken je terug naar het model.

\subsection{Een macht met gekende exponent}

Voorspelt de theorie een verband $y = c \cdot x^k$ waarin je $k$ kent, zoals $y = c \cdot x^2$, stel dan $u = x^k$. Dan is $y = c \cdot u$: een rechte door de oorsprong, als je $y$ uitzet tegen $u$. De kleinste-kwadratenrechte van $y$ tegen $u$ geeft $c$ als richtingscoëfficiënt. Ligt haar snijpunt $b$ dicht bij $0$, dan is dat een extra bevestiging van het model.

\voorbeeld{}{Voor een slinger met lengte $L$ (in m) meet je de periode $T$ (in s). Volgens de natuurkunde is $T^2 = \frac{4\pi^2}{g}L$. Bepaal $g$ uit de metingen.\label{vb:h2-slinger}}{%
\uitleg{Hier is $T^2$ evenredig met $L$. Zet $T^2$ uit tegen $L$: dat moet een rechte door de oorsprong zijn, met richtingscoëfficiënt $\frac{4\pi^2}{g}$.}
\naast{%
{\renewcommand{\arraystretch}{1.2}%
\begin{tabular}{|c|c|c|c|c|c|}
\hline
$L$ & $0{,}2$ & $0{,}4$ & $0{,}6$ & $0{,}8$ & $1{,}0$ \\ \hline
$T$ & $0{,}90$ & $1{,}27$ & $1{,}55$ & $1{,}79$ & $2{,}01$ \\ \hline
$T^2$ & $0{,}81$ & $1{,}61$ & $2{,}40$ & $3{,}20$ & $4{,}04$ \\ \hline
\end{tabular}}
\par\vspace{2mm}
\begin{stappenvrij}{4.2cm}{5.2cm}
T^2 \approx 4{,}025\,L - 0{,}003 & Kleinste-kwadratenrechte van $T^2$ tegen $L$, met $r \approx 1{,}000$. \\
\frac{4\pi^2}{g} \approx 4{,}025 & Het snijpunt is bijna $0$, zoals het model voorspelt. \\
g \approx \frac{4\pi^2}{4{,}025} \approx 9{,}81 & In m/s$^2$.
\end{stappenvrij}}{%
\begin{tikzpicture}[x=3.6cm,y=.62cm,font=\scriptsize]
  \draw[black!14,very thin,xstep=.2,ystep=1] (0,0) grid (1,4);
  \draw[->,grijs] (0,0) -- (1.1,0) node[right,text=black] {$L$};
  \draw[->,grijs] (0,0) -- (0,4.5) node[above,text=black] {$T^2$};
  \foreach \t/\l in {.2/0{,}2,.4/0{,}4,.6/0{,}6,.8/0{,}8,1/1}{\node[below,inner sep=1.5pt] at (\t,0) {$\l$};}
  \foreach \t in {1,2,3,4}{\node[left,inner sep=1.5pt] at (0,\t) {$\t$};}
  \draw[accent,thick] (0,-.003) -- (1.05,4.223);
  \foreach \x/\y in {.2/.81,.4/1.61,.6/2.40,.8/3.20,1/4.04}{\fill[accent] (\x,\y) circle (2pt);}
\end{tikzpicture}}}

\begin{oefeningen}
\oef{h2-kwadraat}
\end{oefeningen}

\subsection{Een exponentieel verband}

Bij een exponentieel verband $y = c \cdot g^x$ staat de veranderlijke in de exponent, en helpt een substitutie zoals $u = x^2$ niet. De logaritme lost dat op: ze maakt van een product een som en van een macht een product.

\begin{theorie}
\begin{stelling}{Lineariseren van een exponentieel verband}
Voor $c > 0$ en $g > 0$ geldt \qquad $\formule{y = c \cdot g^x \iff \ln y = \ln c + x \ln g}$

\vspace{1mm}
Zet je $\ln y$ uit tegen $x$, dan krijg je een rechte met richtingscoëfficiënt $\ln g$ en snijpunt $\ln c$.
\end{stelling}

\begin{bewijs}
Is $y = c \cdot g^x$, dan is $y > 0$, en met de rekenregels van de logaritme is
\bewijsregel{\ln y = \ln c + \ln\bigl(g^x\bigr) = \ln c + x \ln g .}
Omgekeerd volgt uit $\ln y = \ln c + x \ln g$ dat
\bewijsregel{y = \e^{\ln c + x\ln g} = \e^{\ln c} \cdot \bigl(\e^{\ln g}\bigr)^x = c \cdot g^x .}
\end{bewijs}
\end{theorie}

Uit de rechte $\ln y = ax + b$ reken je dus terug met $g = \e^a$ en $c = \e^b$. Het omgekeerde deel van de stelling zorgt ervoor dat dat mag.

\voorbeeld{}{Kroos bedekt een vijver. De bedekte oppervlakte $A$ (in cm$^2$) is na $t$ dagen: $50$; $66$; $88$; $115$; $152$ voor $t = 0, 1, 2, 3, 4$. Stel een exponentieel model op.\label{vb:h2-exponentieel}}{%
\uitleg{De verhoudingen zijn ongeveer constant ($1{,}32$; $1{,}33$; $1{,}31$; $1{,}32$), dus we vermoeden $A = c \cdot g^t$. We zetten $\ln A$ uit tegen $t$, en bepalen de kleinste-kwadratenrechte.}
\begin{stappenbreed}
\ln A\colon\ 3{,}912;\ 4{,}190;\ 4{,}477;\ 4{,}745;\ 5{,}024 & Bijna gelijke stappen van ongeveer $0{,}28$: een rechte. \\
\sum (t_i - \overline{t})^2 = 10,\quad \sum (t_i - \overline{t})(\ln A_i - \overline{\ln A}) \approx 2{,}779 & Met $\overline{t} = 2$ en $\overline{\ln A} \approx 4{,}470$. \\
\ln A \approx 0{,}278\,t + 3{,}914 & $a \approx 0{,}278$, $b \approx 4{,}470 - 0{,}278 \cdot 2$; $r \approx 1{,}000$. \\
g = \e^{0{,}278} \approx 1{,}32,\quad c = \e^{3{,}914} \approx 50 & Terugrekenen met de stelling.
\end{stappenbreed}
\figuren{%
\begin{tikzpicture}[x=.8cm,y=.012cm,font=\scriptsize]
  \draw[black!14,very thin,xstep=1,ystep=40] (0,0) grid (4,160);
  \draw[->,grijs] (0,0) -- (4.5,0) node[right,text=black] {$t$};
  \draw[->,grijs] (0,0) -- (0,180) node[above,text=black] {$A$};
  \foreach \t in {1,...,4}{\node[below,inner sep=1.5pt] at (\t,0) {$\t$};}
  \foreach \t in {40,80,120,160}{\node[left,inner sep=1.5pt] at (0,\t) {$\t$};}
  \draw[accent,thick] plot[domain=0:4.2,samples=40,smooth] (\x,{50*1.32^\x});
  \foreach \x/\y in {0/50,1/66,2/88,3/115,4/152}{\fill[accent] (\x,\y) circle (2pt);}
\end{tikzpicture}\hspace{1.4cm}
\begin{tikzpicture}[x=.8cm,y=1.1cm,font=\scriptsize]
  \draw[black!14,very thin,xstep=1,ystep=.5] (0,0) grid (4,2);
  \draw[->,grijs] (0,0) -- (4.5,0) node[right,text=black] {$t$};
  \draw[->,grijs] (0,0) -- (0,2.3) node[above,text=black] {$\ln A$};
  \foreach \t in {1,...,4}{\node[below,inner sep=1.5pt] at (\t,0) {$\t$};}
  \foreach \t/\l in {0/3{,}5,1/4{,}5,2/5{,}5}{\node[left,inner sep=1.5pt] at (0,\t) {$\l$};}
  \draw[accent,thick] (0,.414) -- (4.2,{.414+.278*4.2});
  \foreach \x/\y in {0/3.912,1/4.190,2/4.477,3/4.745,4/5.024}{\fill[accent] (\x,\y-3.5) circle (2pt);}
\end{tikzpicture}}
Het model is $A \approx 50 \cdot 1{,}32^t$: elke dag komt er ongeveer $32$\,\% bij. Links de gemeten punten met het model, rechts dezelfde punten na de transformatie: een rechte.}

\begin{oefeningen}
\oef{h2-transformeren}
\end{oefeningen}

\subsection{Een macht met onbekende exponent}

Soms weet je dat het verband een macht $y = c \cdot x^k$ is, maar ken je de exponent $k$ niet. Dan kun je niet $y$ tegen $x^k$ uitzetten. Ook hier helpt de logaritme, nu aan beide kanten.

\begin{theorie}
\begin{stelling}{Lineariseren van een machtsverband}
Voor $x > 0$ en $c > 0$ geldt \qquad $\formule{y = c \cdot x^k \iff \ln y = \ln c + k \ln x}$

\vspace{1mm}
Zet je $\ln y$ uit tegen $\ln x$, dan krijg je een rechte met richtingscoëfficiënt $k$ en snijpunt $\ln c$.
\end{stelling}

\begin{bewijs}
Is $y = c \cdot x^k$, dan is $y > 0$ en $\ln y = \ln c + \ln\bigl(x^k\bigr) = \ln c + k \ln x$. Omgekeerd volgt uit $\ln y = \ln c + k\ln x$ dat
\bewijsregel{y = \e^{\ln c} \cdot \bigl(\e^{\ln x}\bigr)^k = c \cdot x^k .}
\end{bewijs}
\end{theorie}

Het verschil met het exponentiële geval: daar transformeer je alleen $y$, hier ook $x$. Een rechte in $(x, \ln y)$ wijst op een exponentieel verband, een rechte in $(\ln x, \ln y)$ op een macht.

\voorbeeld{}{Bij een gitaarsnaar meet je voor verschillende lengtes $L$ (in m) de frequentie $f$ (in Hz) van de toon. Je vermoedt $f = c \cdot L^k$. Bepaal $k$ en $c$.\label{vb:h2-macht}}{%
\uitleg{De exponent is onbekend, dus we zetten $\ln f$ uit tegen $\ln L$, en bepalen de kleinste-kwadratenrechte.}
\figuren{\renewcommand{\arraystretch}{1.2}%
\begin{tabular}{|c|c|c|c|c|c|}
\hline
$L$ & $0{,}30$ & $0{,}40$ & $0{,}50$ & $0{,}60$ & $0{,}65$ \\ \hline
$f$ & $660$ & $495$ & $396$ & $330$ & $305$ \\ \hline
$\ln L$ & $-1{,}204$ & $-0{,}916$ & $-0{,}693$ & $-0{,}511$ & $-0{,}431$ \\ \hline
$\ln f$ & $6{,}492$ & $6{,}205$ & $5{,}981$ & $5{,}799$ & $5{,}720$ \\ \hline
\end{tabular}}
\begin{stappenbreed}
\ln f \approx -1{,}00 \ln L + 5{,}289 & Kleinste-kwadratenrechte, met $r \approx -1{,}000$. \\
k \approx -1{,}00,\quad c = \e^{5{,}289} \approx 198 & Terugrekenen met de stelling.
\end{stappenbreed}
Het model is $f \approx \frac{198}{L}$: de frequentie is omgekeerd evenredig met de lengte. Halveer je de snaar, dan verdubbelt de frequentie, en klinkt de toon een octaaf hoger.}

\begin{werkwijze}
Kies wat je uitzet volgens het model: $y$ tegen $x$ (lineair), $y$ tegen $x^k$ (macht met gekende $k$), $\ln y$ tegen $x$ (exponentieel) of $\ln y$ tegen $\ln x$ (macht met onbekende $k$). Bepaal de kleinste-kwadratenrechte van de getransformeerde gegevens, en reken terug: $g = \e^a$ en $c = \e^b$ bij een exponentieel model, $k = a$ en $c = \e^b$ bij een macht.
\end{werkwijze}

\begin{oefeningen}
\oef{h2-macht}
\end{oefeningen}

%========================================================================================
\section{Een model kritisch bekijken}

Een model is een vereenvoudiging. Drie valkuilen komen in bijna elk onderzoek terug: een hoge $r$ als bewijs zien, te ver voorspellen, en een verband als oorzaak lezen.

\subsection{Residuen bekijken}

Een hoge $r$ bewijst geen lineair verband. Ook punten op een licht gebogen kromme kunnen $r = 0{,}98$ geven. Wat de rechte mist, zie je wel in de residuen.

\begin{theorie}
Een \textbf{\textcolor{accent}{residuplot}} zet de residuen $r_i$ uit tegen de $x_i$. Past het model, dan liggen de residuen willekeurig boven en onder $0$, zonder patroon.
\end{theorie}

\voorbeeld{}{Je knipt vierkanten uit karton met zijde $z$ (in cm) en weegt ze ($m$ in g). Voor $z = 4, 6, 8, 10, 12, 14$ is $m = 0{,}8$; $1{,}9$; $3{,}1$; $5{,}0$; $7{,}3$; $9{,}7$. De kleinste-kwadratenrechte is $m \approx 0{,}894\,z - 3{,}42$, met $r \approx 0{,}987$. Past een lineair model?\label{vb:h2-residuplot}}{%
\naast{%
\uitleg{De $r$ is hoog, maar dat zegt niet genoeg. We berekenen de residuen en kijken of ze een patroon vormen.}
\begin{stappenvrij}{4.4cm}{4.1cm}
\text{residuen } r_i & Meting min voorspelling van de rechte: zie de residuplot. \\
+,\ -,\ -,\ -,\ -,\ + & De tekens: aan de randen positief, in het midden negatief. Een boog. \\
m \text{ tegen } z^2\colon\ r \approx 0{,}9997 & Het karton weegt volgens zijn oppervlakte $z^2$. Nu zijn de residuen hoogstens $0{,}12$, zonder patroon.
\end{stappenvrij}}{%
\begin{tikzpicture}[x=.3cm,y=1.5cm,font=\scriptsize]
  \draw[black!14,very thin,xstep=2,ystep=.25] (2,-.75) grid (16,.75);
  \draw[->,grijs] (2,0) -- (17,0) node[right,text=black] {$z$};
  \draw[->,grijs] (2,-.75) -- (2,1) node[above,text=black] {$r_i$};
  \foreach \t in {4,8,12}{\node[below,inner sep=1.5pt] at (\t,-.75) {$\t$};}
  \foreach \t/\l in {-.5/{-0{,}5},.5/{0{,}5}}{\node[left,inner sep=1.5pt] at (2,\t) {$\l$};}
  \foreach \x/\y in {4/.638,6/-.050,8/-.639,10/-.528,12/-.016,14/.595}{\fill[accent] (\x,\y) circle (2pt);}
\end{tikzpicture}}
Het lineaire model mist een kromming. Het model $m = c \cdot z^2$ past wel, met $c \approx 0{,}050$ g/cm$^2$, en het volgt ook uit de theorie.}

\begin{oefeningen}
\needspace{12\baselineskip}
\oef{h2-residuplot}
\end{oefeningen}

\subsection{Interpoleren en extrapoleren}

Een model is getest op de metingen, en nergens anders. Een rechte die past tussen $x = 1$ en $x = 5$, hoeft niet te kloppen voor $x = 50$. Een kind dat elk jaar $6$ cm groeit, is op $40$ jaar geen $3$ meter.

\begin{theorie}
Een voorspelling met een model binnen het gebied van de metingen heet \textbf{\textcolor{accent}{interpoleren}}. Een voorspelling erbuiten heet \textbf{\textcolor{accent}{extrapoleren}}.
\end{theorie}

\voorbeeld{}{Een zonnebloem is in week $2$ tot $7$ na het zaaien $30$; $46$; $63$; $78$; $95$; $111$ cm hoog. De kleinste-kwadratenrechte is $h = 16{,}2\,w - 2{,}4$, met $r \approx 1{,}000$. Hoe hoog is de plant na $4{,}5$ weken, en na $30$ weken?\label{vb:h2-extrapoleren}}{%
\naast{%
\uitleg{We vullen in, en kijken telkens of de week binnen het gebied van de metingen ligt.}
\begin{stappenvrij}{4.5cm}{4.1cm}
h(4{,}5) = 72{,}9 - 2{,}4 = 70{,}5 & Interpoleren: week $4{,}5$ ligt tussen $2$ en $7$. Betrouwbaar. \\
h(30) = 486 - 2{,}4 \approx 484 & Extrapoleren: bijna $5$ meter. Zonnebloemen stoppen veel vroeger met groeien. \\
h(0) = -2{,}4 & Ook terug in de tijd faalt het model: een negatieve hoogte.
\end{stappenvrij}}{%
\begin{tikzpicture}[x=.36cm,y=.016cm,font=\scriptsize]
  \fill[accent!10] (2,0) rectangle (7,200);
  \draw[black!14,very thin,xstep=2,ystep=50] (0,0) grid (12,200);
  \draw[->,grijs] (0,0) -- (12.8,0) node[right,text=black] {$w$};
  \draw[->,grijs] (0,0) -- (0,220) node[above,text=black] {$h$};
  \foreach \t in {2,4,...,12}{\node[below,inner sep=1.5pt] at (\t,0) {$\t$};}
  \foreach \t in {50,100,150,200}{\node[left,inner sep=1.5pt] at (0,\t) {$\t$};}
  \draw[accent,thick] (2,30) -- (7,111);
  \draw[accent,thick,dashed] (7,111) -- (12.5,200.1);
  \foreach \x/\y in {2/30,3/46,4/63,5/78,6/95,7/111}{\fill[accent] (\x,\y) circle (2pt);}
\end{tikzpicture}}
In de figuur is het gebied van de metingen gekleurd. Daarbuiten is de rechte gestreept: daar weet je niet of ze nog klopt.}

\begin{oefeningen}
\oef{h2-inter}
\end{oefeningen}

\needspace{22\baselineskip}
\subsection{Correlatie en oorzaak}

Twee grootheden die samen stijgen, hebben een hoge $r$. Maar dat betekent niet dat de ene de andere veroorzaakt. Vaak hangen ze allebei af van een derde grootheid.

\voorbeeld{}{In een badstad noteert men een zomer lang per dag het aantal verkochte ijsjes en het aantal reddingsacties op het strand. De correlatiecoëfficiënt is $0{,}8$. Is ijsjes eten gevaarlijk om te zwemmen?\label{vb:h2-oorzaak}}{%
\uitleg{Een hoge $r$ zegt alleen dat de twee samen stijgen en dalen. We zoeken wat ze allebei kan doen stijgen.}
\begin{stappenvrij}{5.6cm}{9.1cm}
\text{derde grootheid: het weer} & Op warme dagen komen er meer mensen naar het strand: ze kopen meer ijsjes, en er zwemmen er meer. \\
\text{nagaan: dagen met hetzelfde weer} & Vergelijk alleen dagen met ongeveer dezelfde temperatuur. Verdwijnt het verband dan, dan kwam het van het weer. \\
\text{oorzaak tonen: een experiment} & Alleen als je zelf één variabele verandert en de rest constant houdt (zie \ref{sec:h2-variabelen}), kun je een oorzaak aantonen.
\end{stappenvrij}
Het besluit \textit{ijsjes zijn gevaarlijk} volgt dus niet uit de gegevens.}

\begin{oefeningen}
\oef{h2-kritisch}
\end{oefeningen}

\subsection{De stappen van het modelleren}

De voorbeelden van dit hoofdstuk volgen telkens dezelfde weg. Samen vormen ze de werkwijze van een onderzoek met metingen.

\begin{werkwijze}
\textbf{1.} Beschrijf de situatie, de variabelen en je aannames. \textbf{2.} Meet zorgvuldig, herhaal waar het kan, en let op systematische fouten. \textbf{3.} Kies een model (lineair, macht, exponentieel) op basis van de theorie of van het spreidingsdiagram. \textbf{4.} Transformeer zo nodig tot een rechte, en bepaal de parameters met de kleinste-kwadratenrechte. \textbf{5.} Controleer het model: liggen de residuen willekeurig rond nul, of zie je een patroon? \textbf{6.} Rapporteer het model met zijn grenzen: het gebied van de metingen, en wat het niet zegt over oorzaken.
\end{werkwijze}

\opdrachtenkop

\begin{oefeningen}
\needspace{14\baselineskip}
\oef{h2-o7}

\oef{h2-o2}

\oef{h2-o8}

\oef{h2-o9}

\oef{h2-o10}

\oef{h2-o11}

\oef{h2-o1}

\oef{h2-o12}

\needspace{14\baselineskip}
\oef{h2-o13}

\oef{h2-o17}

\needspace{14\baselineskip}
\oef{h2-o14}

\oef{h2-o3}

\needspace{14\baselineskip}
\oef{h2-o15}

\oef{h2-o18}

\needspace{12\baselineskip}
\oef{h2-o16}

\oef{h2-o4}

\needspace{14\baselineskip}
\oef{h2-o19}

\oef{h2-o6}

\oef{h2-o21}

\oef{h2-o20}

\oef{h2-o5}
\end{oefeningen}

%----------------------------------------------------------------------------------------
% KOPPEL: verbanden met andere leerstof
%----------------------------------------------------------------------------------------
\needspace{20\baselineskip}
\opdrachtsectie{Koppel}

\lkr{Limieten en afgeleiden, Statistiek en Ruimtemeetkunde: kijk na wat al gezien is. Het pakket legt de volgorde niet vast.}

\begin{oefeningen}
\oef{h2-koppel-1}

\oef{h2-koppel-2}

\oef{h2-koppel-3}
\end{oefeningen}

%----------------------------------------------------------------------------------------
\opdrachtsectie{Test jezelf}

\begin{oefeningen}
\oef{h2-t1}

\oef{h2-t2}

\oef{h2-t3}

\oef{h2-t4}

\oef{h2-t5}

\oef{h2-t6}
\end{oefeningen}

\herhaalschema{%
  \fv{oef:h2-t1}{vb:h2-rechte}{oef:h2-helling}%
  \fv{oef:h2-t2}{vb:h2-exponentieel}{oef:h2-transformeren}%
  \fv{oef:h2-t3}{vb:h2-oorzaak}{oef:h2-kritisch}%
  \fv{oef:h2-t4}{vb:h2-slinger}{oef:h2-kwadraat}%
  \fv{oef:h2-t5}{vb:h2-residu}{oef:h2-residu}%
  \fv{oef:h2-t6}{vb:h2-extrapoleren}{oef:h2-inter}}

\ifzelfstudie
\section*{Maak je eigen samenvatting}
Sluit dit boek en schrijf in je schrift: de soorten variabelen en meetfouten, wat een residu en de kwadratensom zijn, de stellingen over de kleinste-kwadratenrechte (met hun bewijs), de correlatiecoëfficiënt en waarom ze tussen $-1$ en $1$ ligt, hoe je een macht- of exponentieel verband lineair maakt, de stappen van het modelleren, en de drie valkuilen. Vergelijk daarna met de kaders.
\fi
